\chapter{The Collectivization of the Peasant}

Acute anxieties over the grain crisis in the spring of 1929 were veiled in complacent professions of faith in the future. On the basis of the spring sowings a good harvest was predicted. A larger yield was promised from Kolkhozy and Sovkhozy; and a higher proportion of the crop would be marketed. A new procedure was introduced for the grain collections. High quotas were fixed in advance for deliveries from different regions to the collecting agencies. Quotas were assigned to districts and to villages; and within the village pressure was placed on the kulak to bear the main burden of the quota. When the harvest of 1929 was in progress, brigades of party officials, party members, workers and trade unionists were despatched from Moscow, Leningrad and provincial centres to supervise and stimulate the collections. The number of persons engaged in these operations can only be guessed. But the territory was vast, and estimates ranging from 100,000 to 200,000 are not implausible. The peasants---not merely kulaks, but any peasant who had grain that might be thought of as surplus to his own needs---reacted to the campaign by elaborate measures of concealment and by frantic efforts to sell on the free market. Concealment was a criminal offence, and the line between ``trade'', which was legal, and ``speculation'', which was not, was blurred. Reprisals were widely and arbitrarily applied. Failure to fill the quotas was in itself a punishable offence. Kulaks, and alleged kulaks, were fined, sentenced to imprisonment, or simply evicted from the villages. Scenes of violence and brutality occurred. By these means the quotas were filled, and sometimes over-filled. But these results were achieved in conditions of open hostility between the authorities and the peasants, between town and country. Poor peasants were sometimes said to have applauded measures taken against the kulaks. But for the most part solidarity among the peasants prevailed, and kulaks and poor peasants were in collusion to thwart the collections. Party expectations of fanning class war in the countryside were disappointed. 
% 注：去除了段落两端杂乱的连字符和斜杠符号（如 "(regions", "/districts", "f I itself", ".(measures", ".'ll expectations", "f'l disappointed."）。

It was in these unpropitious conditions that the demand for the collective organization of agriculture was insistently pressed, no longer as a distant prospect, but as a solution of current difficulties. The tractor had long been seen as the key to collectivization. In the autumn of 1927 the large Shevchenko Sovkhoz in the Ukraine managed to acquire 60 or 70 tractors, which were organized in ``tractor columns'' to work its own fields and those of neighbouring Kolkhozy or peasant holdings. The example was imitated elsewhere; and in 1928 Shevchenko established the first Machine Tractor Station (MTS) with a park of tractors to be leased out to Kolkhozy and Sovkhozy in the region. In June 1929 a central office, Traktortsentr, was set up in Moscow to organize and control a network of state MTSs. Peasant prejudices against the innovation, and perhaps against the degree of state intervention involved in it, were hard to overcome. Tractors were sometimes denounced as the work of Anti-Christ. The success of the experiment seemed, however, to have been limited mainly by the supply of tractors; in the autumn of 1929 only 35,000, most of them of American manufacture, were available for the whole of the USSR. Everywhere it came, the tractor was a powerful agent of collectivization. 
% 注：去除了 ".iiTu been" 等断行乱码。

The revival of the Kolkhozy which started in 1927 had at first led to a proliferation of small and loosely organized Kolkhozy, whose performance was unsatisfactory. In the middle of 1928 a drive began for ``large'' Kolkhozy, defined as Kolkhozy with a sown area of 2000 hectares, sufficiently large to be worked by tractors. But at this moment the Kolkhozy were outstripped by the Sovkhozy. At the party central committee in July 1928, Stalin called for the creation of large-scale grain-growing Sovkhozy, which were thought of as ``grain factories'' working on industrial lines. The prototype of the new Sovkhozy was one appropriately named Gigant, which covered 41,000 hectares of hitherto largely uncultivated land in the North Caucasian region; and this was followed by similar enterprises in the Volga, Ural and Siberian regions. The tractor and the MTS were indispensable pre-requisites of this operation, which was later sometimes criticized as ``gigantomania''. When collectivization began in earnest, enthusiasm for the Sovkhozy waned, and they were once more eclipsed by the Kolkhozy. 
% 注：去除了 "I' I The", "-f of as" 等边缘字符。

One problem keenly debated in party circles was the question what to do with the kulak, or the peasant labelled as such by the authorities, the peasant who commonly farmed the largest and best plots of land in the village, was best equipped with animals and machines, produced and held the largest surpluses of grain, and offered the strongest opposition to Soviet policies, including the policy of collectivization. Opinions were sharply divided. If the kulak, together with his land and inventory, were incorporated in the Kolkhoz, he would---so some party members argued---make an important contribution to its production and efficiency. But he would also---as others reasonably predicted---exercise a dominating influence over it, and guide it in directions hostile to the purposes of the party and the state. If, however, he were excluded from the Kolkhoz, what was to become of him? He could not be allowed to retain his land and possessions, and constitute an independent unit of production side by side with the Kolkhoz. He would have to be evicted and expelled from the region; and this was a harsh measure which few at first were ready to contemplate. No acceptable solution could be found. 

Throughout the summer and autumn of 1929 the drive at the centre for more and more collectivization increased in intensity. But two assumptions continued to be made even by its most enthusiastic promoters. The first, whatever pressure might be put on the peasant by local authorities, was that collectivization would be voluntary; the second, whatever emphasis might be placed on the urgency of the operation, was that it would take, at any rate, some years to complete. By the end of the year, the leaders had talked themselves out of both these assumptions, and were suddenly prepared to take the plunge into forced and immediate collectivization of Soviet agriculture as a whole. The decisive change seems to have been brought about by two pre-disposing factors. The first was a mood of desperation bred by the annual nightmare of the grain collections; apart from the prospect of increased production, Kolkhozy could be compelled, more easily than the individual peasant, to deliver their grain to the official agencies. The second was a mood of elation inspired by the successes of industrialization and by the prospects of the five-year plan. Agriculture was after all a form of industry. If the forcing of the rate of industrialization had fulfilled the hopes even of the most optimistic, it would show a lack of faith to reject the promise of a forced rate of collectivization. Unflinching resolution was the quality required to carry the position by assault. 
% 注：去除了 "'agencies", "th^^", "industry^", "Ilack" 等标点和乱码，修正了 "thie" 为 "the"，"industn^zation" 为 "industrialization"。

Stalin, as was usual with him, refused to enter the arena till the issue was clarified by debate, and the moment was ripe for decision. From April to November 1929 he remained silent. Then he published in Pravda the customary article on the anniversary of the revolution under the heading ``The Year of the Great Break-Through''. Having hymned the triumphs of industrialization and the development of heavy industry, Stalin turned to agriculture, which had achieved ``a fundamental break-through from small backward individual economy to large-scale progressive collective agriculture''. The middle peasant, he claimed, ``has entered the Kolkhozy''. The kulak was barely mentioned. As regards the future: 
% 注：去除了 "The 'kulak" 中的单侧引号，增加了上一句末尾遗漏的引号。

\begin{quote}
If the development of Kolkhozy and Sovkhozy proceeds at an accelerated pace, there are no grounds for doubting that in three years or so our country will become a great grain country, if not the greatest in the world. 
\end{quote}

The article contained a vision of the future and an analysis of the present, but no call for immediate action. Considering its character as a celebratory pronouncement, it was restrained and cautious. The party still hovered on the brink of a decision which it hesitated to take. 
% 注：去除了 "^_j^.|jfThe", "^ I of the" 的非常严重的乱码。

At the session of the party central committee a few days later the tone was already sharper. Stalin repaired the omission in his article by referring to a ``mass offensive of poor and middle peasants against the kulak''. Molotov was the most uncompromising of several speakers who sought to force the pace of collectivization. He rejected the estimate of the five-year plan (which had modestly foreseen the collectivization of 20 per cent of the sown area in the five-year period) as too extended a projection; most regions would be totally collectivized by 1931, some by the autumn of 1930. The kulak was denounced as an ``undefeated enemy'', who should not be allowed to penetrate the Kolkhozy. But none of the speeches were published till collectivization was already under way, so that their increasingly urgent tone was not known to the party or to the public. The resolutions adopted at the end of the session were less precise than Molotov's pronouncement about dates, perhaps reflecting scepticism among some members of the committee, but called for ``a decisive offensive against the kulak, containing and cutting off attempts of kulaks to penetrate the Kolkhozy''. The question what to do with the kulak was still not faced. During the next few weeks enthusiastic reports poured in from party organs in the principal grain-growing regions on the progress of collectivization; and on December 5, 1929, the Politburo appointed a commission with instructions to submit, within two weeks, a draft decree on the rate of collectivization in various regions. The commission included representatives of the regions, but no member of the Politburo, and was evidently thought of as a technical, not a policy-making, body. 
% 注：去除了 "i", "fnembers", ".gainst", "o penetrate" 等断行乱码，并修正 "Stahn" 为 "Stalin"。

The fragmentary records of the commission published many years later may reflect the confusion of the proceedings. The commission broke up into sub-commissions which canvassed many bold proposals; one of them seems to have coined the phrase ``the liquidation of the kulaks as a class''. The draft submitted to the Politburo by the commission on December 22 was, however, still relatively cautious. It proposed the collectivization of the principal grain-growing regions in from two to three years (stipulating that in some districts and regions progress might be more rapid), and of other regions in from three to four years; a warning was added against ``an ecstasy of dictation''. It was assumed that kulaks could not be admitted to the Kolkhozy; their means of production, i.e. their machines and their animals, were to be transferred to the Kolkhozy, and distant and inferior land was to be allocated to them. Recalcitrant kulaks were to be expelled from the region; those who submitted might be allowed to work in some undefined capacity in the Kolkhozy. 
% 注：去除了 "■^1/K", "gainst", "jcould", "/(from", "n I work" 等极其杂乱的边缘符号。

Before the Politburo could consider the commission's report, a conference of Agrarian Marxists assembled in Moscow; Stalin seized the occasion to deliver to it his first public speech for many months. It was the most violent attack yet made on the kulak. ``Dekulakization'' or ``the liquidation of the kulaks as a class'' was described as ``one of the most decisive turns in our whole policy''. About the same time, an active party worker of Kalmyk origin named Ryskulov, who had been a member of the Politburo commission, criticized its report in a note addressed to the Politburo. It was an odd gesture, which could hardly have been risked without higher approval. Ryskulov called for an increase in the pace of collectivization, and its extension to cotton-growing and live-stock regions, which had been overlooked in the draft, as well as for the surrender to the Kolkhozy of animals, including milch-cows and poultry, which the draft had proposed to leave in the possession of the individual peasant. The draft was revised in the light of these observations, and the revised text adopted by the party central committee on January 5, 1930. 

The resolution of January 5, 1930 was the key decision in the process of collectivization. It proclaimed ``the replacement of large kulak production by large Kolkhoz production'', and ``the liquidation of the kulaks as a class''. Collectivization of the principal grain-growing regions---the Lower and Middle Volga and the North Caucasus---was ``perhaps in the main'' to be completed by the autumn of 1930 or the spring of 1931, and of the other grain-growing regions in the autumn of 1931 or the spring of 1932. The supply of tractors and machines was to be expedited, but this was not to be treated as a condition of collectivization. An awkwardly worded paragraph provided that, in the transitional period, ``the fundamental means of production (animals and implements, farm buildings, live-stock reared for sale)'' should be vested in agricultural cooperatives within the Kolkhoz. The fate of the kulaks---evidently still a controversial issue---was not yet settled. A further commission was set up under Molotov to deal with it, and on January 30, 1930, the Politburo adopted a resolution. Its text has never been published, but its substance was sufficiently indicated by its title ``On Measures to Eliminate Kulak Households in Districts of Total Collectivization''. 
% 注：去除了 "■in", "[of", "las", "fneans", "ive-stock", ":ooperatives" 的边缘符号。

What happened in the countryside in the winter of 1929--1930 was determined not so much by the texts of resolutions as by the character of the operation which was mounted to carry them out. During the winter 25,000 industrial workers (said to have been selected from 70,000 volunteers) were assigned to permanent work in rural areas. These were merely the nucleus of a large army of party activists, officials, agricultural experts, tractor mechanics and Red Army men, dispersed over the countryside to shepherd the peasant into the new Kolkhozy. Considerable attention was paid to organization; military terms like ``brigade'', ``headquarters'' and ``staff'' were in use. All concerned received elaborate briefings. In some places courses of instruction were set up for peasants. But few of those responsible had any experience of the countryside, or of peasant life and mentality. The instructions themselves were confused and contradictory; and excess of zeal in interpreting them seemed a venial fault. The proclaimed intention not to apply compulsion to middle or poor peasants was soon frustrated. Since no mercy could be shown to the kulak, who was treated as an enemy of the regime, any peasant resisting collectivization was liable to be branded as a kulak, or, as being hand-in-glove with the kulaks, subjected to the same penalties. Tens of thousands of kulaks were evicted from their holdings and dwellings, and turned adrift to fend for themselves or deported to remote regions; their animals, machines and implements were turned over to the Kolkhoz. Few peasants of any category entered the Kolkhoz voluntarily. What the peasants most resented was the demand to surrender their animals. Many chose to slaughter these rather than hand them over. Throughout the campaign the line between persuasion and compulsion was thinly drawn. 
% 注：跨页导致的断裂 "winter of 1929--1930 was determined..." 已修复。去除了 "he nucleus", "ural experts", "or themselves" 中的断字和符号。

A feature of the operation was the demand for larger and larger Kolkhozy---an extension of the ``gigantomania'' which had begun with the Sovkhozy. Giant Kolkhozy were formed in the principal grain-growing regions, of which the largest covered 80,000 hectares. But the main purpose of the Kolkhozy distinguished them from the Sovkhozy, being not so much to bring virgin lands into cultivation as to combine existing small Kolkhozy and peasant holdings into large units. These Kolkhozy, which might include several villages and several thousand peasant households, were steps, consciously taken, on the road to collectivization, meaning that the whole of the land in a given area was included in one or more comprehensive Kolkhozy; such localities were described as ``areas of total collectivization''. Much publicity was given to a petition from the Khoper district in the Lower Volga region to become an area of total collectivization, the process to be completed within the period of the first five-year plan. This was hailed as a model. But two major obstacles impeded the expansion of the Kolkhozy---their unpopularity among the majority of the peasants, who clung tenaciously to the possession of their own plots of land and animals, and the inadequate supply of tractors and other machines, which alone gave meaning and purpose to the policy of collectivization. A further grave handicap was a deficiency of personnel---both of party members and Soviet officials having any contact with the countryside or knowledge of its problems, and of the agronomists, veterinary workers and skilled mechanics essential to the working of so vast a transformation. 
% 注：去除了 "’not", "[Jlhe", "Ijwith", "/the", "I essential", "Vast" 等错位字符。修正 "dehciency" 为 "deficiency"。

The widespread confusion resulting from these proceedings, and sporadic disorders among the peasantry, threatened the spring sowings and frightened the authorities. An article by Stalin, published on March 2, 1930, under the title ``Dizzy from Success'', called a halt to further collectivization. Pressure was relaxed; and during the spring many peasants who had been pressed into the Kolkhozy were allowed to leave them. The retention of small individual holdings and of some animals was now again tolerated. The retreat seems to have been signalled in time to allow sowings to proceed more or less normally. This fortunate step, combined with exceptionally favourable weather, accounted for the record grain harvest of 1930, the largest since the revolution. But a sharp decline in the number of animals was ominous for the future; and the respite was short. The hammer-blows of the first months of the year had broken the back of peasant resistance, and shattered beyond repair the old peasant order. The kulak had been expelled or crushed. In areas of total collectivization, the mir was formally abolished by a decree of July 30, 1930. When the collectivization drive was resumed at the end of the year, it encountered less active opposition, and proceeded more rapidly. By the middle of 1931 two-thirds of all holdings in the main grain-growing regions had been incorporated in the Kolkhozy, and the rest followed during the next few years. 
% 注：去除了 "(land", "//spring", "if 1930", "n the", "Ijjshattered", "/been", "Imir", "^d of", "Vnore", "Holdings", "jcorporated" 等边缘严重错乱的字符。

The full costs of the transformation were, however, not long delayed. Production had been disorganized. The most efficient producers had been thrown out. Though the supply of tractors and machinery slowly increased, the Kolkhozy were not yet equipped to fill the gap. What were more efficient were the grain collections; a higher proportion of the harvest was extracted from the Kolkhozy than had been obtained from individual peasants. The peasant went hungry. More and more animals were slaughtered because they could no longer be fed. Bad harvests both in 1931 and in 1932 crowned the calamity. Grain was still remorselessly collected even from the worst-hit areas; and in the ensuing winter what had been the richest grain-growing regions were prey to a famine which was worse than anything experienced eleven years earlier after the civil war (see p. 36 above). The number of those who died of hunger cannot be computed. Estimates have varied from one to five million. 

Collectivization completed the agrarian revolution which had begun in 1917 with the seizure of landlords' estates by the peasants, but which had left unchanged the ancient methods of cultivation and peasant ways of life. The final stage, unlike the first, owed nothing to spontaneous peasant revolt; Stalin aptly called it ``revolution from above'', but wrongly added that it had been ``supported from below''. In the past twelve years agriculture had remained a quasi-independent enclave in the economy, functioning on its own lines, and resisting any attempt from without to alter them. This was the essence of NEP. It was an uneasy compromise which did not last. Once a powerful central authority in Moscow had taken in hand the planning and reorganization of the economy, and embarked on the path of industrialization, and once the failure of agriculture under the existing system to supply the needs of a rapidly expanding urban and factory population became evident, the break logically followed. The battle was joined, and was fought out with great tenacity and bitterness on both sides. 

The ambition of the planners was to apply to agriculture the two great principles of industrialization and modernization. The Sovkhozy were conceived as mechanized grain factories. The mass of peasants were to be organized in Kolkhozy constituted on the same model. But extravagant hopes of ensuring a supply of tractors and other machines sufficient to make such a project viable in practical terms were disappointed. The party had never had any firm foothold in the countryside. Neither the leaders who took the decisions in Moscow, nor the army of party members and supporters who descended on the countryside to enforce them, had any understanding of peasant mentality, or any sympathy for the ancient traditions and ancient superstitions which were the core of peasant resistance. Mutual incomprehension was complete. The peasant saw the emissaries from Moscow as invaders who had come, not only to destroy his cherished way of life, but to reconstitute the conditions of slavery from which the first stage of the revolution had freed him. Force was on the side of the authorities, and was brutally and ruthlessly applied. The peasant---and not only the kulak---was the victim of what looked like naked aggression. What was planned as a great achievement ended in one of the great tragedies that left a stain on Soviet history. The tiller of the soil had been collectivized. But it took Soviet agriculture many years to recover from the disaster which accompanied the process. It was not till the later nineteen-thirties that the production of grain returned to the levels achieved before forced collectivization began; and the loss in the number of animals persisted still longer. 
% 注：修复了由于分栏和跨页造成的极其严重的符号噪点。去除了 "\ constituted", "lensuring", "l/to make", "f/pointed.", "hrm foothold", "de.scended", "fife", "-9 production", "■V of" 等等。修正 "fife" 为 "life"。